\section{Formalism}

\subsection{Definition of the series}
Label the levels with the Latin letters \(a, b, \dots, z\) --- \nlevels{} levels in all ---
and write the value of the \(n\)-th level as \(L(n)\), so that \(L(1)=a\) and
\(L(\nlevels)=z\). The series is fixed by a starting value and a recurrence:
\[
  a = L(1) = \avalue, \qquad L(n+1) = L(n)^{\,L(n)}.
\]
This is (offset) \emph{tetration}: the base of the exponent itself grows at every step.

\subsection{The crypto coma}
By the \emph{crypto coma} we mean the self-power of the last level:
\[
  \boxed{\;\ccoma = z^{z} = L(\nlevels)^{\,L(\nlevels)}.\;}
\]

\subsection{Compact form: the self-power operator}
The whole ladder is one operation applied many times. Introduce the \emph{self-power
operator} \(\star\):
\[
  x^{\star} = x^{x} \qquad (\text{one star} = \text{``raise the number to itself''}).
\]
Writing \(\star^{k}\) for \(k\) applications, the whole ladder and its summit fit on one
line:
\[
  a = \avalue, \qquad L(n) = \bigl(\avalue\bigr)^{\star\,(n-1)}, \qquad
  \boxed{\;\ccoma = \bigl(\avalue\bigr)^{\star\,\nlevels}.\;}
\]
In other words: \(\avalue\), raised to itself \nlevels{} times --- one star per letter of
the alphabet (\(a\) has no stars, \(z\) has 25, and the crypto coma \(=z^{\star}\) has
\nlevels).

\subsection{Estimating the magnitude (on a log scale)}
Already \(b = a^{a} = \bigl(\avalue\bigr)^{\avalue} = \bvalue\); this number cannot be
written in digits. So we work with the \emph{order of magnitude} \(M(n)=\log_{10}L(n)\).
Taking \(\log_{10}\) of the recurrence gives a convenient formula for the order:
\[
  M(1) = \amagnitude, \qquad M(n+1) = M(n)\cdot 10^{\,M(n)}.
\]
Thus \(M(2) = \bmagnitude\), and from level \(c\) on the order itself no longer fits in a
machine number and is represented as a \emph{tower} of powers of ten. The crypto coma
\(\ccoma\) is a base-ten power tower about \ccheight{} storeys tall.

\subsection{The symbol}
For the crypto coma we propose the symbol \(\ccoma\) --- a handwritten ``C'' (Crypto Coma)
struck through by a double upward arrow \(\uparrow\uparrow\), evoking the rising tower of
powers (and Knuth's notation for tetration).

\subsection{``Cardinality''}
In the spirit of set theory we appraise the crypto coma as a \emph{cardinality} and submit
it to the attention of the scientific community and everyone who cares about numbers.
Strictly speaking, \(\ccoma\) is a finite natural number; on the subtlety of this usage see
Section~\ref{sec:disclaimer}.
